\section{Constructing the character field and its arithmetic host}
\label{sec:arithmetic-hosts}

\status{AH1--AH18 reconstruct the finite character construction, its actual
cyclotomic fixed field, every field of \texttt{HostDer}, the half-different
ideal and its trace-self-dual inverse, and the literal \texttt{fieldChi}
conclusion. The hypotheses on the primes $q,d$, their splitting congruences,
and their eventual analytic windows remain hypotheses. These arguments
construct arithmetic data once those hypotheses hold; they do not select
the analytic windows.}

Throughout this section $U_a=(\mathbb Z/a\mathbb Z)^\times$ for $a>0$,
$T_e=2^{e+2}$, and $H_e=U_{T_e}/\langle-1\rangle$.
The symbol $\varphi(a)$ denotes Euler's totient, whereas $\phi$ denotes a
frame. Write $n=\phi.\mathrm{ord}=|A|$, $n_0=|A_0|$ and
$o=n/n_0$. We use multiplicative notation for the finite group $A$;
transporting it to an additive group on $\{0,\ldots,n-1\}$ gives the
literal frame. Put
\[
 G=\{\pm1\}\times\{\pm1\}\times A,\qquad
 c=(-1,-1,1),\quad \iota_d=(1,-1,1),\quad
 \Psi=\{(s,t,a):t=1\}.                         \tag{AH.1}
\]
The exact prime hypothesis is
\[
 q,d\text{ prime},\qquad q\equiv3\pmod8,
 \quad d\equiv7\pmod{16},\qquad
 \left(\frac{-d}{q}\right)=1.                 \tag{AH.2}
\]
It implies that $q,d$ are distinct odd primes. In particular no choice of
such primes is included in the arguments below.

\paragraph{AH1: finite character duality and the two-adic quotient.}
The quotient $H_e$ is cyclic of order $2^e$. Indeed,
$v_2(5^{2^j}-1)=j+2$: the case $j=0$ is immediate, and the induction step
uses $5^{2^j}+1\equiv2\pmod8$ for $j\ge1$, with the first step checked
directly. Thus $5$ has order $2^e$ modulo $T_e$. Its powers are $1$ modulo
$4$, so $-1$ is outside this subgroup, and counting the $2^{e+1}$ odd
residue classes proves $U_{T_e}=\langle-1\rangle\times\langle5\rangle$.

Let $M>0$ and let $\Xi$ be a subgroup of the complex characters of $U_M$.
Set
\[
 B=U_M/\operatorname{ann}(\Xi),\qquad
 \operatorname{ann}(\Xi)=\{u:\chi(u)=1\text{ for every }\chi\in\Xi\}.
                                                               \tag{AH.3}
\]
Then $|B|=|\Xi|$, and the characters of $B$ are exactly the characters in
$\Xi$. To see both assertions, sum $\chi(u)$ over $u\in U_M$ and
$\chi\in\Xi$. The sum of a nontrivial character is zero, since translating
the sum by a point where the character is not $1$ multiplies it by that
value. Summing first over $\Xi$ gives
$|\operatorname{ann}(\Xi)||\Xi|=|U_M|$. Every character in $\Xi$ factors
through $B$, and a finite abelian group has exactly its order many complex
characters: on each cyclic factor of order $a$ one chooses any $a$th root
of unity. The resulting cardinality equality proves the assertion.
This argument includes groups with odd-order factors.

We also use the following elementary reduction fact. If $a\mid b$,
the map $U_b\to U_a$ is onto. Lift a unit modulo $a$ to an integer and
use CRT to make it $1$ at every prime dividing $b$ but not $a$.
The resulting integer is coprime to $b$ and has the required image.

\paragraph{AH2: the minimal modulus and the actual mixed frame.}
Fix $e\ge0$, put $W_e=\operatorname{lcm}(T_e,M)$, and let
$\xi_W:U_{W_e}\to B$ be reduction modulo $M$ followed by the quotient.
Choose the least positive integer $Q=Q_e$ satisfying
\[
 Q\mid W_e,\quad T_e\mid Q,\quad
 u\equiv1\pmod Q\Longrightarrow\xi_W(u)=1.       \tag{AH.4}
\]
The set is nonempty because $W_e$ belongs to it. Reduction $U_{W_e}\to U_Q$
is surjective; its kernel lies in $\ker\xi_W$, so there is a unique
surjective homomorphism $\xi:U_Q\to B$ with
$\xi(u\bmod Q)=\xi_W(u)$. Define
\[
 \operatorname{cl}:U_Q\longrightarrow H_e,
 \qquad \eta(u)=(\operatorname{cl}(u),\xi(u)),
 \qquad A=\eta(U_Q)\subset H_e\times B.          \tag{AH.5}
\]
The first projection $p_1:A\to H_e$ is onto. On $\ker p_2$ it is
injective, since an element with both coordinates $1$ is the identity.
Hence $\ker p_2$ is cyclic. Choose a generator $a_0$ and set
$A_0=\langle a_0\rangle=\ker p_2$. If $p_1(a_0^j)=1$, then $a_0^j=1$;
this is precisely the generator condition of \texttt{Frame}.
Relabeling the finite group by $\operatorname{Fin}(|A|)$ gives the literal
\texttt{ModData.frame}, with its supplied group isomorphism, not merely a
group of the same order. Let $a:U_Q\to A$ denote $\eta$ with this relabeling,
and let $\xi_A:A\to B$ denote $p_2$. Then
\[
 p_1a=\operatorname{cl},\quad \xi_Aa=\xi,
 \quad A_0=\ker\xi_A,
 \quad o=[A:A_0]=|B|=|\Xi|,
 \quad 2^e\mid n.                              \tag{AH.6}
\]
If every element of $B$ has order dividing $2^j$, then $A$ is a $2$-group,
because it is a subgroup of $H_e\times B$. No such assumption is made
for the construction itself.

\paragraph{AH3: CRT constructs the Artin map, including its sign twist.}
More generally take any $Q>0$, $T_e\mid Q$, and any homomorphism
$\xi:U_Q\to B$, and form $A$ as in AH2. Assume (AH.2),
$\gcd(Q,qd)=1$, $e\ge2$, and
$(-d/\ell)=1$ for every odd prime $\ell\mid Q$.
Set $N=Qqd$ and
\[
 \alpha(u)=\left(\frac{u}{q}\right),\qquad
 \beta(u)=\left(\frac{u}{d}\right),\qquad
 w(s)=a(s\bmod Q),\quad s\in\{\pm1\}.
\]
The literal homomorphism is
\[
 \operatorname{art}(u)=
 \bigl(\alpha(u),\beta(u),a(u\bmod Q)w(\alpha(u))\bigr).
                                                               \tag{AH.7}
\]
For a prescribed $(s,t,b)\in G$, choose $v\in U_Q$ with
$a(v)=b\,w(s)$. Choose a residue of Legendre sign $s$ modulo $q$ and
one of sign $t$ modulo $d$; one may use $1$ or $-1$ because both primes
are $3$ modulo $4$. CRT gives a unit $u\in U_N$ with these three
residues. Since $w(s)^2=1$, (AH.7) maps it to $(s,t,b)$.
Thus the Artin map is surjective.

The remaining identities follow from the same formula:
\[
 \begin{aligned}
 \operatorname{art}(-1)&=c,\\
 p_1(\operatorname{art}(u)_A)&=[u\bmod T_e],\\
 u\equiv1\pmod{Qq}&\Longrightarrow
              \operatorname{art}(u)\in\{1,\iota_d\},\\
 u\equiv1\pmod{Qd}&\Longrightarrow
              \operatorname{art}(u)^2=1.
 \end{aligned}                                                   \tag{AH.8}
\]
For the first identity the $A$ coordinate is $a(-1)w(-1)=1$;
for the second, $p_1w(s)=1$ because $s=\pm1$ dies in $H_e$.
These calculations prove every Artin field of
\texttt{ModData.charData}. The factor $w(\alpha(u))$ is indispensable
when the original characters are not even.

\paragraph{AH4: the cyclotomic fixed field and the actual quadratic roots.}
Take any \texttt{CharData}, with its surjection
$\operatorname{art}:U_N\to G$ and the identities (AH.8), and put
\[
 L=\mathbb Q(\zeta_N),\qquad
 U=\ker(\operatorname{art}),\qquad E=L^U.        \tag{AH.9}
\]
Under the usual cyclotomic identification
$\tau_u(\zeta_N)=\zeta_N^u$, the fundamental theorem of Galois theory gives
\[
 \operatorname{Gal}(E/\mathbb Q)\simeq U_N/U\simeq G,
 \qquad [E:\mathbb Q]=4n.                       \tag{AH.10}
\]
The isomorphism is characterized by
$\operatorname{gal}(\tau_u|_E)=\operatorname{art}(u)$, so the inclusion
$E\hookrightarrow L$ satisfies the literal \texttt{art\_spec}.

For completeness, the quadratic elements can be constructed inside this
same cyclotomic field. For a prime $p\equiv3\pmod4$ dividing $N$, choose
$\zeta_p=\zeta_N^{N/p}$ and set
$g_p=\sum_{a\bmod p}(a/p)\zeta_p^a$. Changing $a$ to $u^{-1}a$ gives
$\tau_u g_p=(u/p)g_p$. Also $\overline{g_p}=-g_p$ and
\[
 g_p\overline{g_p}
 =\sum_{t\ne0}(t/p)\sum_{b\ne0}\zeta_p^{(t-1)b}=p.
\]
Here the inner sum is $p-1$ when $t=1$ and $-1$ otherwise, and
$\sum_{t\ne0}(t/p)=0$. Hence $g_p^2=-p$.
If $(u/p)$ is a character of $\operatorname{art}(u)$, its value is $1$
on $U$, so $g_p\in E$ with the prescribed Galois action. In particular
$s_d=g_d$ realizes the required second sign. For the Artin map of AH3,
$s_q=g_q$ realizes the first sign as well. This proves
\texttt{CharData.toHostRaw}, \texttt{ModData.host} and
\texttt{AnyGrp.host}, including their specified roots, embeddings and
Galois maps.

\paragraph{AH5: exactness of the modulus and the primitive conductor.}
Return to the minimal $Q_e$ of AH2 and assume $e\ge1$.
For every prime $\ell\mid Q$ there is a unit $v\in U_Q$ with
\[
 v\equiv1\pmod{Q/\ell},\qquad \eta(v)\ne1.      \tag{AH.11}
\]
If $T_e\mid Q/\ell$ and no such $v$ existed, every unit modulo $W_e$
congruent to $1$ modulo $Q/\ell$ would have trivial $\xi_W$-value.
Thus $Q/\ell$ would satisfy (AH.4), contradicting minimality.
Otherwise $\ell=2$ and $v_2(Q)=e+2$. The unit $v=1+Q/2$ is $1$ modulo
$Q/2$ but is neither $1$ nor $-1$ modulo $T_e$; the latter assertion uses
$e\ge1$ and $Q/2\equiv2^{e+1}\pmod{T_e}$. Its $H_e$ coordinate is
therefore nontrivial. Lift $v$ by CRT to a unit $u$ which is $1$ modulo
$q$ and $d$. Then $u\equiv1\pmod{N/\ell}$ and
$\operatorname{art}(u)\ne1$. This is every clause of
\texttt{HostRaw.Exact}, for the actual modulus $Q$.

When $\Xi=\langle\psi\rangle$ for a primitive character of conductor $M$,
\[
                    M\mid Q,\qquad Q=W_e.       \tag{AH.12}
\]
Indeed, let $g=\gcd(M,Q)$ and let $v\in U_M$ be $1$ modulo $g$.
Generalized CRT gives $z\equiv v\pmod M$ and $z\equiv1\pmod Q$.
It is a unit modulo $W_e$, because $T_e\mid Q$; (AH.4) implies
$\psi(v)=1$. Thus $\psi$ factors through $U_g$. Primitivity forces
$g=M$, proving $M\mid Q$; together with $T_e\mid Q\mid W_e$ this gives
the displayed equality. The literal source only needs its divisibility
consequence.

\paragraph{AH6: one eventually constant multiplier in the frame order.}
For the fixed $M,\Xi$, define
\[
 J_e=\xi(\ker\operatorname{cl})\subset B.
\]
Projection of $A\subset H_e\times B$ onto $H_e$ has kernel
$\{1\}\times J_e$. Consequently
\[
              n_e=2^e|J_e|,\qquad 1\le|J_e|\le|\Xi|.             \tag{AH.13}
\]
These subgroups form a descending sequence. To check this without assuming
$Q_e\mid Q_{e+1}$, lift a representative at level $Q_{e+1}$ to a unit
modulo $W_{e+1}$, reduce it modulo $W_e$, then modulo $Q_e$.
Its two-adic residue, originally $\pm1$ modulo $T_{e+1}$, is $\pm1$
modulo $T_e$, and its $B$ coordinate is unchanged. Hence
$J_{e+1}\subset J_e$. The nonincreasing positive integers $|J_e|$ attain
a minimum. Choose an index $e_1$ attaining it and write that minimum as
$c_0$. For all $e\ge e_1$,
\[
 1\le c_0\le|\Xi|,\qquad n_e=c_0\,2^e.          \tag{AH.14}
\]
The constants precede $q,d$ and every analytic scale. This is the exact
quantifier order of \texttt{ordAny}.

\paragraph{AH7: the prescribed characters are actually kept.}
For $\chi\in\Xi$, let $\chi_B:B\to\mathbb C^\times$ be its descent and
write $\theta_A=\chi_B\xi_A$. Define the character of the actual group $G$
by
\[
 \theta_\chi(s,t,b)=\theta_A\bigl(b\,w(s)\bigr). \tag{AH.15}
\]
It is trivial on $(1,1,A_0)$ and on $\iota_d$, so belongs to the literal
$X^+$ character set. Substitution in (AH.7) cancels $w(\alpha)^2$ and gives
\[
       \theta_\chi(\operatorname{art}(u))=\chi(u)                 \tag{AH.16}
\]
whenever both sides are read on units of a common multiple of $N$ and $M$.
Thus the primitive Dirichlet character induced by $\theta_\chi$ is the
primitive character induced by $\chi$. This uses uniqueness of the
primitive inducing character, which follows by factoring both characters
through the common modulus and then through the gcd of their conductors.
Accordingly every primitive character in \texttt{Level.ofSub} is kept in
$X^+$.

The first-sign character is not among them. Use CRT to choose
$u\equiv1\pmod Q$, $u\equiv1\pmod d$ and $(u/q)=-1$.
For every character (AH.15), its value on $\operatorname{art}(u)$ is $1$,
whereas the first sign is $-1$. Equal primitive inducing characters would
have equal values on all units modulo $N$, a contradiction. Since
\texttt{Level.ofSub.psK=none}, this is the entire
\texttt{ExceptionalCharacterKept} predicate. It does not assert that all
characters in $X^+$ are kept.

\paragraph{AH8: the exact kernel and every allowed quadratic twist.}
Let $\psi$ be primitive of conductor $M$, and $\Xi=\langle\psi\rangle$.
For every integer $a$ coprime to $N$, (AH.7) and (AH.12) give
\[
 \begin{split}
 \operatorname{art}(a)=1\quad\Longleftrightarrow\quad&
 a\equiv\pm1\pmod{T_e},\quad\psi(a)=1,\\
 & (a/q)=(a/d)=1. 
 \end{split}                                                     \tag{AH.17}
\]
Indeed, the first two coordinates must be $1$; then $w(\alpha)=1$ and
the remaining coordinate is trivial exactly when both components of
$\eta(a)$ are trivial. The annihilator of $\langle\psi\rangle$ is
$\ker\psi$. This proves \texttt{KerChi}, including $M\mid Q$.

Now let $\theta:G\to\mathbb C^\times$ be trivial on $(1,1,A_0)$.
Its restriction to $A$ factors through $A/A_0\simeq B$, hence by AH1 is
some $\chi\in\Xi$. The restrictions to
$s\mapsto(s,1,w(s))$ and $t\mapsto(1,t,1)$ have values $s^i,t^j$
for some $i,j\in\{0,1\}$. This follows because the image of the generator
of a group of order $2$ has square $1$, hence is $1$ or $-1$.
The factorization
\[
 \operatorname{art}(a)
  =(\alpha,1,w(\alpha))(1,\beta,1)(1,1,a(a\bmod Q))
\]
therefore gives
\[
 \theta(\operatorname{art}(a))
   =\chi(a)(a/q)^i(a/d)^j,\qquad
 \theta(\iota_d)=(-1)^j.                        \tag{AH.18}
\]
This is \texttt{TwistForm} with its prescribed $\iota_d$ sign. It includes
all characters trivial on the designated cyclic subgroup, with no
restriction to even characters or to $2$-groups.

\paragraph{AH9: cyclotomic reduction gives the complete local Galois laws.}
Let $L=\mathbb Q(\zeta_N)$ and let $P_L$ lie above a rational prime $p$.
Write $N=p^v m$, with $p\nmid m$. In the residue field, the reduction of
$\zeta_N$ has exact order $m$. Its order divides $m$ because a root of
$p$-power order reduces to $1$. A root of order prime to $p$ keeps its
order: if $\eta\ne1$, $\eta^a=1$, $p\nmid a$, and $\eta\equiv1\pmod{P_L}$,
then $1+\eta+\cdots+\eta^{a-1}=0$ would reduce to $a=0$, a contradiction.
Applying this to powers of $\zeta_N^{p^v}$ gives exactness. Thus
\[
 \zeta_N^r\equiv\zeta_N^s\pmod{P_L}
       \quad\Longleftrightarrow\quad r\equiv s\pmod m.          \tag{AH.19}
\]

Here one may test integral residue maps on $\zeta_N$, since
$\mathcal O_L=\mathbb Z[\zeta_N]$. The relevant integral statement has
a direct local proof. For each $p$, the prime-to-$p$ polynomial
$X^m-1$ has nonzero derivative at all its roots modulo $p$.
Hensel factorization therefore makes $\mathbb Z_p[\zeta_m]$ a product
of unramified discrete valuation rings. On each factor,
$\Phi_{p^v}(1+X)$ is Eisenstein if $v>0$: its constant coefficient is
$p$, its leading coefficient is $1$, and the others are divisible by $p$.
Adjoining $\zeta_{p^v}-1$ consequently gives a discrete valuation ring,
with the power basis of this Eisenstein polynomial. Since
$\mathbb Z_p[\zeta_N]=\mathbb Z_p[\zeta_m,\zeta_{p^v}]$, this is the
full integral closure locally. The finite index between
$\mathbb Z[\zeta_N]$ and $\mathcal O_L$ is therefore divisible by no prime,
so is $1$. For $v=0$ the unramified argument already suffices.

It follows from (AH.19), on the integral generator, that
\[
 \begin{aligned}
 \tau_u\in I(P_L)&\Longleftrightarrow u\equiv1\pmod m,\\
 \tau_u\text{ induces arithmetic Frobenius at }P_L
      &\Longleftrightarrow u\equiv p\pmod m.
 \end{aligned}                                                   \tag{AH.20}
\]
For the second assertion the two ring maps on the residue field agree on
the generator precisely when $\overline{\zeta_N}^u=
\overline{\zeta_N}^{p}$. A unit $u$ in this congruence class exists by
CRT, choosing its $p$-part to be $1$. The decomposition group is the
preimage of the cyclic residue Galois group generated by this Frobenius,
so
\[
 \tau_u\in D(P_L)
       \quad\Longleftrightarrow\quad
       u\equiv p^j\pmod m\text{ for some }j\ge0.                \tag{AH.21}
\]

These formulas descend to every Galois subfield $E$ of $L$ under its
actual restriction map. More explicitly, restriction maps decomposition
and inertia groups onto the respective groups at $P=P_L\cap\mathcal O_E$.
For decomposition, lift an automorphism of $E$ to $L$ and adjust the
lift by an element of $\operatorname{Gal}(L/E)$ so that it fixes $P_L$;
this is possible because that group is transitive on the primes above
$P$. For inertia, the resulting residue action fixes $\kappa(P)$.
The decomposition group over $E$ maps onto
$\operatorname{Gal}(\kappa(P_L)/\kappa(P))$; cancel this residue action
by an element over $E$. This proves the claimed surjectivity, including
the inertia direction that restriction alone would not prove.
The residue surjectivity used here is the usual finite Galois prime
theorem: the stabilizer acts on the residue field with kernel inertia
and quotient the finite-field Galois group. It follows by localizing
the normal extension at the prime and lifting the embeddings of its
maximal unramified subextension.

Consequently, for the cyclotomic realization of any raw host, writing
$N_p=N/p^{v_p(N)}$, one has the literal identities
\[
 \begin{aligned}
 \sigma\in I(P)&\Longleftrightarrow
     \exists u\in U_N: u\equiv1\pmod{N_p},\quad
                       \operatorname{gal}(\sigma)=\operatorname{art}(u),\\
 \sigma\text{ is arithmetic Frobenius at }P&\Longleftrightarrow
     \exists u\in U_N: u\equiv p\pmod{N_p},\quad
                       \operatorname{gal}(\sigma)=\operatorname{art}(u),\\
 \sigma\in D(P)&\Longleftrightarrow
     \exists u\in U_N,\ j\ge0: u\equiv p^j\pmod{N_p},\quad
                       \operatorname{gal}(\sigma)=\operatorname{art}(u).
 \end{aligned}                                                   \tag{AH.22}
\]
For the Frobenius converse, two lifts induce the same residue Frobenius
exactly when their quotient is in inertia. This supplies the necessary
lift without imposing that a general lift itself be Frobenius. If $p\nmid N$,
the inertia group is trivial and the unique arithmetic Frobenius has
image $\operatorname{art}(p\bmod N)$.

\paragraph{AH10: the CM field, the exact CM-type and the sign characters.}
Now let $H$ be an arbitrary \texttt{HostRaw}, not necessarily one just
constructed in AH4. Its specified embedding in $\mathbb Q(\zeta_N)$
and \texttt{art\_spec} imply that $\operatorname{art}$ is onto: extend
an automorphism of $E$ to the cyclotomic field and read its exponent on
$\zeta_N$. They also identify the image of $E$ with the fixed field of
$\ker\operatorname{art}$. Thus AH9 applies to this actual raw host.

Complex conjugation on $\zeta_N$ is $\tau_{-1}$, so its restriction to $E$
is $\operatorname{gal}^{-1}(c)$; denote it by $\overline{\phantom{x}}$.
Every embedding of $E$ into $\mathbb C$ extends to the cyclotomic field,
and hence
$\psi(\overline{x})=\overline{\psi(x)}$ for every embedding $\psi$.
Because $s_d^2=-d$ and $d>0$, no embedding of $E$ is real. The fixed field
$F=E^{\langle\overline{\phantom{x}}\rangle}$ is totally real, is the
maximal real subfield, and $E/F$ is quadratic. Therefore
\[
 [E:\mathbb Q]=4n,\qquad [F:\mathbb Q]=2n,\qquad
 x\in F\Longleftrightarrow\overline{x}=x.        \tag{AH.23}
\]
An automorphism induces conjugation at even one complex embedding exactly
when it is this automorphism, by injectivity of that embedding. This
proves the full \texttt{conj\_iff} predicate.

The actual CM-type consists of the embeddings with
$\psi(s_d)=i\sqrt d$. Each embedding sends $s_d$ to one of
$\pm i\sqrt d$, and conjugation exchanges the two possibilities.
Thus the type has $2n$ elements, selects exactly one from each conjugate
pair, and, for $\psi$ in the type,
\[
 \psi\sigma\text{ is in the type}
       \quad\Longleftrightarrow\quad
       \operatorname{gal}(\sigma)\in\Psi.        \tag{AH.24}
\]
Also $K=\mathbb Q(s_d)$ is exactly the fixed field of $\Psi$.

There is a nonzero first-sign eigenvector even when a raw host does not
assert \texttt{PsQ}. Take the quadratic fixed field of the kernel of the
first-sign character, choose $x$ moved by conjugation there, and use
$a=x-\overline{x}\ne0$. Automorphisms of first sign $1$ fix it and those
of first sign $-1$ negate it. This proves \texttt{ps\_root}; it does not
replace the stronger square-root-of-$-q$ supplied in AH4 for our constructed
host.

Finally compare the image of the raw element $s_d$ with the Gauss sum
$g_d$ of AH4. Their squares are both $-d$, so they differ by sign. The
Gauss action therefore forces
\[
 \operatorname{art}(u)_d=(u/d).                  \tag{AH.25}
\]
For $p\nmid N$, AH9, (AH.25), quadratic reciprocity and $d\equiv3\pmod4$
give every coordinate required by \texttt{frob}: the first sign is the
defined $\mathrm{psAt}(p)$, the second is $(p/d)=(-d/p)$, and the $H_e$
coordinate is the class of $p$ modulo $T_e$.

\paragraph{AH11: ramified decomposition and the unramified CM extension.}
The quadratic integral calculation in AC8 gives
\[
 \mathcal O_K=\mathbb Z[\omega],\quad
 \omega=(1+s_d)/2,\quad
 f_\omega(X)=X^2-X+(d+1)/4,quad \operatorname{disc}(K)=-d.
                                                               \tag{AH.26}
\]
At $2$ this polynomial is $X(X-1)$ modulo $2$, since $d\equiv7\pmod{16}$.
At an odd prime different from $d$ its roots are distinct, and it splits
exactly when $(-d/p)=1$; at $d$ the prime ramifies. Thus every prime
dividing $Qq$ splits in $K$. Restriction of a decomposition automorphism
at such a prime fixes its prime in $K$. In a split quadratic extension
that stabilizer is trivial, so
\[
              p\mid Qq\quad\Longrightarrow\quad D(P)\subset\Psi.
                                                               \tag{AH.27}
\]
This proves \texttt{decomp\_typeI}, including the prime $2$.

At $d$, (AH.22) and $N_d=Qq$ give
\[
                     I(P)=\{1,\iota_d\}.        \tag{AH.28}
\]
The containment follows from \texttt{art\_d}. For the reverse inclusion,
CRT gives $u\equiv1\pmod{Qq}$ and $u\equiv-1\pmod d$.
Its second sign is $-1$ by (AH.25), so its image is $\iota_d$.

We need a further first-sign fact at $d$. The integer $-d$ is a square
modulo $Qq$. At every odd prime factor, start with the prescribed
nonzero square root and lift by solving the linear congruence for its
next digit; $2s$ is invertible. At $2$, start with $s=1$ modulo $8$.
If $s^2\equiv-d\pmod{2^a}$, $a\ge3$, replacing $s$ by
$s+2^{a-1}$ changes its square by $2^a$ modulo $2^{a+1}$, and therefore
corrects the next bit when necessary. CRT combines the prime powers.
The first sign of $\operatorname{art}$ factors through $U_{Qq}$ by
\texttt{art\_d}. If $s^2\equiv-d\pmod{Qq}$, then $d\equiv-s^2$, whence
\[
 u\equiv d\pmod{Qq}\quad\Longrightarrow\quad
              \operatorname{art}(u)_q=-1.       \tag{AH.29}
\]
This proof uses the first sign of $\operatorname{art}(-1)$ and the square
of a sign; it does not assume \texttt{PsQ} for a general raw host.

Conjugation fixes no prime above a divisor of $N$. For a divisor of $Qq$
this follows from (AH.27), since $c\notin\Psi$. At $d$, choose a
Frobenius lift $F$ with exponent $u\equiv d\pmod{Qq}$. If conjugation
were in the decomposition group, then
$c=\rho\,\operatorname{gal}(F)^j$ with
$\rho\in\{1,\iota_d\}$ and $j\ge0$. The $A$ coordinate and
\texttt{art\_h} would give $d^j\equiv\pm1\pmod{T_e}$.
Reducing modulo $16$, the congruence $d\equiv7$ forces $j$ even, because
$7^j$ is $7$ for odd $j$ and $1$ for even $j$.
But the first sign, using (AH.29), gives $-1=(-1)^j$, a contradiction.
Thus \texttt{decompConj} holds. At primes outside $N$ inertia is trivial;
at primes dividing $N$ conjugation is not even in the decomposition group.
It follows that conjugation is in no finite inertia group, proving the
literal \texttt{unramEF}: $E/F$ is unramified at every finite prime.

\paragraph{AH12: splitting, ordinary places and the exact $\eta$ clauses.}
The factorization in (AH.26) gives
\[
 \#\{\mathfrak p\subset\mathcal O_K:\mathfrak p\mid p\}=2
 \quad\Longleftrightarrow\quad
 p=2\ \text{or}\ \bigl(p\ne2\text{ and }(-d/p)=1\bigr).        \tag{AH.30}
\]
For a prime $P$ of $E$, the source's ordinary condition says
$\sigma\tau\in\Psi\Longleftrightarrow\sigma\in\Psi$ for every
$\tau\in D(P)$ and every $\sigma$. Taking $\sigma=1$ shows this is
equivalent to $D(P)\subset\Psi$; the converse follows because $\Psi$ is
a subgroup. Since restriction maps $D(P)$ onto the decomposition group
in $K$, this is equivalent to splitting in $K$. This proves the full
\texttt{ordinary} equivalence for all rational primes, including $d$.

Since $E/F$ is quadratic and unramified at finite primes, every nonzero
prime of $F$ either has two primes above it or has one prime of residue
degree $2$. The latter case occurs exactly when conjugation fixes a prime
$P$ above it. For $p\nmid N$, the decomposition group is generated by
the arithmetic Frobenius $\sigma$, so this is equivalent to
$c\in\langle\operatorname{gal}(\sigma)\rangle$.
For $p\mid N$, AH11 shows it never occurs. With $\eta_F=1$ in the split
case and $-1$ in the inert case, these are precisely all three clauses of
\texttt{HostDer.eta}.

\paragraph{AH13: constructing the actual half-different ideal.}
Let $\mathfrak D=\mathfrak D_{E/K}$.
The relative inertia group for $E/K$ is $I(P)\cap\Psi$. At a prime above
$d$ it is trivial by (AH.28); outside $N$ it is trivial by AH9.
In a finite extension of number fields, a prime divides the different
exactly when it is ramified. This follows locally from the trace dual
over discrete valuation rings, or from the derivative of a local integral
generator; in the unramified case that derivative is a unit.
Therefore
\[
 P\mid\mathfrak D\quad\Longrightarrow\quad P\mid Qq.           \tag{AH.31}
\]
Every such rational prime splits in $K$. Hence, for
$\mathfrak p=P\cap\mathcal O_K$,
$\overline{\mathfrak p}\ne\mathfrak p$.

The different is invariant under every $\sigma\in\operatorname{Gal}(E/\mathbb Q)$.
Indeed, $K$ is stable because the Galois group is abelian, and
\[
 \operatorname{Tr}_{E/K}(\sigma z)
       =\sigma|_K\bigl(\operatorname{Tr}_{E/K}(z)\bigr).
\]
This preserves integrality in $\mathcal O_K$, hence preserves the
codifferent and its inverse $\mathfrak D$. In particular the exponents
of $P$ and $\overline P$ in $\mathfrak D$ agree.
Order the finite set of primes of $K$ below the support of $\mathfrak D$,
and define
\[
 \mathfrak b=
 \prod_{\substack{P\mid\mathfrak D\\
              P\cap\mathcal O_K<\overline P\cap\mathcal O_K}}
                       P^{v_P(\mathfrak D)}.                  \tag{AH.32}
\]
Exactly one prime of each conjugate pair contributes. Unique factorization
of nonzero ideals gives
\[
 \mathfrak D=\mathfrak b\overline{\mathfrak b},\qquad
 \mathfrak b+\overline{\mathfrak b}=\mathcal O_E.                 \tag{AH.33}
\]
An element of $\operatorname{Gal}(E/K)=\Psi$ preserves contractions to
$K$, commutes with conjugation, and preserves the exponents of the
different. It therefore preserves the selected product:
$\sigma\mathfrak b=\mathfrak b$ for every $\sigma\in\Psi$.
This proves every clause of the source theorem named
\texttt{KcDiff.principal}. Its conclusion is a selected integral ideal
$\mathfrak b$. It does not assert that this ideal is generated by one
element and uses no class-number hypothesis.

\paragraph{AH14: the inverse ideal is self-dual for the conjugated trace.}
The codifferent is
\[
 \mathfrak D^{-1}
   =\{z\in E:\operatorname{Tr}_{E/K}(z\mathcal O_E)
                                      \subset\mathcal O_K\}.
\]
For any nonzero fractional $\mathcal O_E$-ideal $J$, its trace dual satisfies
\[
 J^*=\{x:\operatorname{Tr}_{E/K}(xJ)\subset\mathcal O_K\}
       =\mathfrak D^{-1}J^{-1}.                  \tag{AH.34}
\]
To prove the equality, the trace condition is equivalent to
$xJ\subset\mathfrak D^{-1}$: for $y\in J$ and $a\in\mathcal O_E$,
the element $ay$ is still in $J$. Since a nonzero ideal in a Dedekind domain
is invertible, the latter inclusion is equivalent to
$x\in\mathfrak D^{-1}J^{-1}$.

Take any integral ideal $\mathfrak b$ satisfying (AH.33), not necessarily
the particular choice (AH.32). It is nonzero because the different is
nonzero. Conjugation maps $\mathfrak b^{-1}$ bijectively onto
$(\overline{\mathfrak b})^{-1}$: multiply the defining inclusion
$y\mathfrak b\subset\mathcal O_E$ by conjugation in both directions.
Apply (AH.34) with $J=(\overline{\mathfrak b})^{-1}$ to obtain
\[
 J^*=(\mathfrak b\overline{\mathfrak b})^{-1}
             \overline{\mathfrak b}=\mathfrak b^{-1}.
\]
Consequently, for every $x\in E$,
\[
 x\in\mathfrak b^{-1}
 \quad\Longleftrightarrow\quad
 \operatorname{Tr}_{E/K}(x\overline y)\in\mathcal O_K
                 \text{ for every }y\in\mathfrak b^{-1}.        \tag{AH.35}
\]
This proves \texttt{KcDiffSmulMem.unimodular} in both directions, with
the literal $\mathcal O_K$-valued trace pairing and the literal inverse
fractional ideal. No principal generator or metric comparison has been
substituted for this pairing.

\paragraph{AH15: a finite-kernel formula for the root discriminant.}
Let $E\subset L=\mathbb Q(\zeta_N)$ be any subfield, and let
$U\subset U_N$ be its fixing subgroup. For $f\mid N$ define
\[
 K_f=\ker(U_N\longrightarrow U_f),\qquad
 \rho_U(f)=\frac{|K_f\setminus U|}{|K_f|},\qquad
 S(U)=\sum_{a\mid N}\Lambda(a)\rho_U(N/a),        \tag{AH.36}
\]
where $\Lambda(p^j)=\log p$ for $j\ge1$ and is zero on other positive
integers. Then
\[
           \log\operatorname{rd}(E)=S(U),\qquad
 \operatorname{rd}(E)=|\operatorname{disc}E|^{1/[E:\mathbb Q]}.
                                                               \tag{AH.37}
\]
Here is a derivation, so this formula introduces no analytic input.
By AH9, $\mathcal O_L=\mathbb Z[\zeta_N]=\mathcal O_E[\zeta_N]$.
If $f_E$ is the minimal polynomial of $\zeta_N$ over $E$, the relative
different is $(f_E'(\zeta_N))$. The monogenic different formula follows
directly from Lagrange interpolation: for a monic separable polynomial
of degree $a$, the sums
$\sum_{f_E(z)=0}z^j/f_E'(z)$ are $0$ for $j<a-1$ and $1$ for $j=a-1$.
Reducing higher powers by $f_E$, these identities give a trace-dual basis
of $\mathcal O_E[\zeta_N]$ in
$f_E'(\zeta_N)^{-1}\mathcal O_E[\zeta_N]$, with a triangular change of
basis of determinant $\pm1$. This proves the ideal formula in both
directions. Also
\[
 f_E'(\zeta_N)=
       \prod_{\substack{u\in U\\u\ne1}}(\zeta_N-\zeta_N^u).
\]
The discriminant tower identity, obtained from the determinant of the
trace pairing in a tower, is
\[
 |\operatorname{disc}L|
   =|\operatorname{disc}E|^{[L:E]}
      N_{L/\mathbb Q}(\mathfrak D_{L/E}).
\]
Subtract its logarithm from the corresponding product for $E=\mathbb Q$.
This gives
\[
 [L:E]\log|\operatorname{disc}E|
       =\sum_{u\notin U}\log|N_{L/\mathbb Q}(\zeta_N-\zeta_N^u)|.
                                                               \tag{AH.38}
\]
If $h$ is the additive order of $u-1$ modulo $N$, then $h>1$ and
$\zeta_N-\zeta_N^u$ is a root of unity times $1-\eta$ for a primitive
$h$th root $\eta$. Its norm has logarithm
$\varphi(N)\Lambda(h)/\varphi(h)$. Indeed,
$\Phi_h(1)=p$ for $h=p^j$ and $1$ otherwise. The last evaluation follows
by induction from
$\prod_{a\mid h,\,a>1}\Phi_a(1)=h$, obtained by evaluating
$(X^h-1)/(X-1)$ at $1$. Taking the norm first in $\mathbb Q(\eta)$ and
then in $L$ proves the stated logarithm.

For each divisor $h>1$ of $N$, the elementary identity
\[
 \frac{\Lambda(h)}{\varphi(h)}
   =\sum_{\substack{a\mid N\\h\mid a}}
             \Lambda(a)\frac{\varphi(N/a)}{\varphi(N)}           \tag{AH.39}
\]
holds. If $h$ is not a prime power both sides vanish. If $h=p^j$ and
$N=p^v m$, $p\nmid m$, the right side is the geometric totient sum
$\log p\,\sum_{i=j}^v\varphi(p^{v-i})/\varphi(p^v)
=\log p/\varphi(p^j)$.
The condition $h\mid a$ says $u\equiv1\pmod{N/a}$, and
$|K_f|=\varphi(N)/\varphi(f)$ by surjective reduction.
Interchange the two finite sums in (AH.38), use (AH.39), and divide by
$[L:E][E:\mathbb Q]=\varphi(N)$. This proves (AH.37).

\paragraph{AH16: all three numerical discriminant bounds and its support.}
First, for any subgroup $U\subset U_N$,
\[
                  [U_N:U]\le\exp(2S(U)).        \tag{AH.40}
\]
Choose the least divisor $c_1$ of $N$ for which $K_{c_1}\subset U$.
If $K_a,K_b\subset U$, then $K_{\gcd(a,b)}\subset U$:
for $u\equiv1\pmod{\gcd(a,b)}$, CRT gives a unit $v$ with
$v\equiv1\pmod a$, $v\equiv u\pmod b$, and then
$u=v(v^{-1}u)\in K_aK_b$. Lift from the lcm modulus to $N$ using AH1.
It follows by minimality that $c_1\mid f$ whenever $K_f\subset U$.
If $c_1\nmid f$, the proper subgroup $K_f\cap U$ has index at least $2$
in $K_f$, so $\rho_U(f)\ge1/2$. Counting prime powers shows
\[
 \sum_{\substack{a\mid N\\c_1\nmid N/a}}\Lambda(a)=\log c_1:
\]
for $p^v\Vert N$ and $p^b\Vert c_1$ the sum contains precisely the last
$b$ powers of $p$. Hence $2S(U)\ge\log c_1$, while
$[U_N:U]\le\varphi(c_1)\le c_1$. This proves (AH.40).

Apply this to $U=\ker\operatorname{art}$ of a raw host.
Surjectivity and $|G|=4n$ give
\[
                         4n\le\operatorname{rd}(E)^2.           \tag{AH.41}
\]
For the upper bound, $\rho_U(f)\le1-\varphi(f)/\varphi(N)$ because
$1\in K_f\cap U$. At $f=N/d=Qq$, the image of $K_f$ has order $2$ by
(AH.28), so $\rho_U(f)=1/2$. At $f=N/q=Qd$, the group $K_f$ is cyclic,
being isomorphic to $U_q$ by CRT, and \texttt{art\_q} makes its image
have exponent dividing $2$. Thus its image has order at most $2$, and
$\rho_U(N/q)\le1/2$. In the remaining terms of (AH.36), retain just the
improvement at powers of $2$ over the bound $\rho\le1$.
Writing $N=2^v m$ with $m$ odd, their total improvement is
\[
 \log2\sum_{j=1}^v\frac{\varphi(N/2^j)}{\varphi(N)}
       =\log2.
\]
Since $q,d$ are distinct odd primes, these corrections are disjoint.
Using $\sum_{a\mid N}\Lambda(a)=\log N$ gives
\[
 \begin{split}
 \log\operatorname{rd}(E)
 &\le\log N-\log2-\tfrac12\log q-\tfrac12\log d,\\
 \operatorname{rd}(E)&\le\frac Q2\sqrt{qd}.       \tag{AH.42}
 \end{split}
\]

For the lower bound, \texttt{art\_h} implies that every $u\in U$ reduces
to $1$ or $-1$ modulo $T_e$. Write again $N=2^v m$, $m$ odd; here
$v\ge e+2$. For $0\le i\le e+1$, the reduction of $K_{2^i m}$ in
$U_{T_e}$ has order $2^{e+1}/\varphi(2^i)$ and equal-size fibers.
Among its elements, at most two of $\{1,-1\}$ are possible if $i=0,1$,
and at most one if $i\ge2$. Therefore
\[
 \rho_U(2^i m)\ge
     1-\frac{\varphi(2^i)+\mathbf1_{i<2}}{2^{e+1}}.
\]
Use in (AH.36) the distinct prime-power terms $a=2^{v-i}$ for these
indices and the term $a=d$. All other terms are nonnegative.
Since $\sum_{i=0}^{e+1}\varphi(2^i)=2^{e+1}$ and
$\rho_U(N/d)=1/2$, this proves
\[
 \begin{split}
 \log\operatorname{rd}(E)
   &\ge(e+1-2^{-e})\log2+\tfrac12\log d,\\
 2^{e+1-2^{-e}}\sqrt d&\le\operatorname{rd}(E).  \tag{AH.43}
 \end{split}
\]
The counts also justify the stated reductions when $i=0$, since
$U_1$ is the trivial group.

No rational prime outside $N$ divides $\operatorname{disc}E$, by AH9 and
the ramification criterion for the different. The prime $d$ does divide
it because its inertia group is nontrivial. The same is true of $2$:
CRT makes $K_m$ surject onto $U_{T_e}$, and \texttt{art\_h} makes its
image surject onto the nontrivial group $H_e$; thus $I_2\ne1$.
Finally $E/F$ is unramified at every finite prime, so the discriminant
tower identity gives $|\operatorname{disc}E|=|\operatorname{disc}F|^2$.
There are $2n$ conjugate pairs of complex embeddings of $E$, an even
number, so its discriminant has positive sign. Thus
\[
 \operatorname{disc}E=(\operatorname{disc}F)^2,\qquad
 2d\mid\operatorname{disc}E,\qquad
 p\mid\operatorname{disc}E\Longrightarrow p\mid N
 \quad(p\text{ prime}).                         \tag{AH.44}
\]
These are the literal \texttt{discr\_sq}, \texttt{discr\_primes},
\texttt{rd\_ge}, \texttt{rd\_le} and \texttt{rd\_sq\_ge} fields.

\paragraph{AH17: the full raw-to-derived host theorem.}
The preceding arguments prove \texttt{HostRaw.der} for every raw host,
with no new arithmetic premise. The exact coverage of its thirty fields
is as follows:
\[
\begin{array}{p{0.72\linewidth}l}
\texttt{art\_surj, cm, conj\_iff, finrank\_E, finrank\_F, mem\_F,
cmType, ps\_root, art\_ch}&\mathrm{AH10}\\[2pt]
\texttt{frob, frob\_any, inertia, decomp}&\mathrm{AH9,AH10}\\[2pt]
\texttt{ps\_d, inertia\_d, decomp\_typeI, decompConj, unramEF}
 &\mathrm{AH11}\\[2pt]
\texttt{splitK, ordinary, eta}&\mathrm{AH12}\\[2pt]
\texttt{diffSupport, principal}&\mathrm{AH13}\\[2pt]
\texttt{unimodular}&\mathrm{AH14}\\[2pt]
\texttt{discr\_sq, discr\_primes, rd\_ge, rd\_le, rd\_sq\_ge}
 &\mathrm{AH15,AH16}\\[2pt]
\texttt{norm\_mu}&\mathrm{AC2,AC3}.
\end{array}                                                     \tag{AH.45}
\]
For the last line, the raw hypotheses already give
$[E:\mathbb Q]=4n$, $d>3$ prime, and $K=\mathbb Q(\sqrt{-d})$;
AC2--AC3 prove $N_{E/K}(\zeta)=1$ for every root of unity, including its
$2$-primary part. Thus that result introduces no dependency on
\texttt{HostDer.norm\_mu} itself. Appending these proved fields to $H$
gives the literal \texttt{HostAx.ofRaw H}, whose raw projection is $H$.
In particular the arithmetic input to HC4's host witness is now supplied.

\paragraph{AH18: the complete prescribed-character field theorem.}
Let $\psi$ be any primitive complex Dirichlet character of conductor $M>0$,
and set $\Xi=\langle\psi\rangle$, $o=|\Xi|$. Choose $c_0,e_1$ by AH6
and set
\[
                   e_0=\max(e_1,2,4M,64o).       \tag{AH.46}
\]
For every $e\ge e_0$, every $q,d$ satisfying (AH.2), and every integer
$b\ge4M$, assume additionally
\[
 q\ne3,\qquad \gcd(M,qd)=1,\qquad
 (-d/p)=1\text{ for every odd prime }p\le b.     \tag{AH.47}
\]
Use the frame from AH2, the Artin map from AH3, the fixed field from AH4,
and AH17 to obtain an actual $H:\texttt{HostAx}\,\phi$.
Then $\phi.e=e$, $n=c_0\,2^e$, $H.q=q$, $H.d=d$, and
$\phi.o=o=(\texttt{Level.cyc}\,\psi).o$.

Here are all the remaining quantitative rows. Since
$Q\mid W_e\mid4M2^e$ and $n\ge2^e$,
\[
 Q\mid\operatorname{lcm}(T_e,M),\qquad Q\le4Mn.
\]
Also $4M\le e<2^e$ implies
$Q\le4M2^e\le2^{2e}\le n^2$. The bound $o\le2^o$ gives
$o^{64}\le2^{64o}\le2^e$. If an odd prime $\ell$ divides $Q$, it divides
$M$; hence $4\ell\le4M\le b$, and the splitting hypothesis in (AH.47)
applies. Thus the eight fields of \texttt{FieldRows} are all proved:
the first-sign Legendre formula and the actual square root of $-q$ come
from AH3--AH4, exactness from AH5, $Q\le n^2$ and $o^{64}\le2^e$ from
these inequalities, $q\ne3$ and splitting below $b$ from (AH.47), and
$4\ell\le b$ for odd $\ell\mid Q$ from the divisor argument.
AH7 supplies \texttt{KeptInXp} and
\texttt{ExceptionalCharacterKept}; AH8 supplies \texttt{KerChi} and
\texttt{TwistForm}. This is every conjunct of the literal
\texttt{UpperTheory.Top.fieldChi} theorem, with $c_0\ge1$ and with
$c_0,e_0$ chosen before $e,q,d,b$.

The result does not select primes satisfying the simultaneous analytic
pair window, does not supply \texttt{VVle}, and does not prove a top-window
estimate. The original $q,d$ congruence and splitting hypotheses are
preserved exactly. HC10 can now use this constructed field, and SZ13 can
construct ranks over its actual frame when the required order window is
supplied. Those joins retain their analytic inputs and the separate
\texttt{HostGood}, \texttt{WinOK}, \texttt{LevelPrime},
$d_u$--$D_s$, \texttt{Over} and \texttt{RecOK} interfaces.

